% =====================================================================
% BioPAWS-2 section — standalone fragment for the merged manuscript.
% Intended insertion point: as a new \subsection right after
% \subsection{Training} and before \subsection{Benchmark performance},
% OR as a dedicated \section{BioPAWS-2: an instruction-tuning benchmark}
% immediately after the Introduction. Labels: sec:biopaws2.
% Numbers are real measurements (RTX PRO 6000 Blackwell, 2026-06).
% =====================================================================

\section{BioPAWS-2: a unified instruction-tuning benchmark for bio-foundation models}
\label{sec:biopaws2}

\subsection{Motivation: from a probe to a benchmark}

The original BioPAWS \citep{wang2026biopaws} was introduced as a \emph{probe} of a single
scientific claim --- that general-purpose language syntax transfers to biological
sequence structure --- and contained five pairwise / single-sequence tasks. As a
\emph{benchmark} it has two limitations that we address here. First, it is narrow: it does
not cover the breadth of biological sequence classification, regression, and cross-modal
tasks against which a foundation model should be measured. Second, and more fundamentally,
the entire established evaluation landscape for biological sequence models --- TAPE
\citep{rao2019tape}, PEER \citep{xu2022peer}, FLIP \citep{dallago2021flip}, ProteinGym
\citep{notin2023proteingym}, and the genomic GUE suite \citep{zhou2023dnabert2} --- is
built around the \emph{protein-language-model-plus-classification-head} paradigm: each task
ships as (sequence, label) rows, and a model is evaluated by freezing a backbone, attaching
a task-specific head, and training that head. This paradigm cannot evaluate a generative
chat model in its native interface, and conversely a benchmark of free-form biological
question answering (e.g. BixBench \citep{bixbench2024}) cannot be answered by a specialized
PLM. There is no common axis on which a specialist PLM and a generalist LLM can be compared.

BioPAWS-2 closes this gap by re-expressing biological sequence tasks as
\emph{instruction-tuning question answering}. Every item is a chat record: an instruction
that names the task and enumerates the candidate labels, an input carrying the sequence(s),
and a target answer that is one of the enumerated labels (for classification), an ordinal
bucket (\texttt{low}/\texttt{medium}/\texttt{high}, for regression, with the raw value
retained for a Spearman back-channel), or free text (for description and reasoning tasks).
This single format (i) lets a generalist LLM answer zero-shot or be fine-tuned directly,
(ii) lets a specialist PLM run its native head-training protocol on the same train/test
split, and (iii) plugs directly into downstream agent systems, where a uniform chat
interface is the integration point. Crucially, BioPAWS-2 is a public \emph{training
resource}, not only a test set: the SFT corpus is itself the contribution, and the model
evaluated in this paper (OmniGene-4-MM) is merely one of many models that can be trained
and scored on it.

\subsection{Construction and coverage}

BioPAWS-2 organizes tasks into nine families spanning the breadth of biological sequence
analysis: F1 pairwise alignment/homology, F2 protein single-sequence function, F3 DNA /
genomic classification, F4 variant-effect prediction, F5 structure-as-text, F6 cross-modal
(DNA + protein + small molecule), F7 biological chain-of-thought reasoning, F8 biomedical
sequence-grounded QA, and F9 multimodal (image + text) QA. The label/answer conversion
follows four rules: native classification becomes multiple-choice QA with the candidate
set in the instruction; regression becomes ordinal-bucket QA; pairwise tasks become binary
comparison QA; and description/captioning tasks remain free-form (scored by ROUGE-L/BLEU).
The instruction template explicitly lists all candidate labels (``\texttt{The result will
be one of the following: A, B, C}''), which generalizes from two-class to $k$-class tasks
and forces label-grounded answering. To prevent prompt overfitting, each task draws from a
small bank of paraphrased templates.

The current release instantiates \textbf{306K QA examples across all nine families} (22
tasks) with real data, combining: (i) the BioPAWS protein- and DNA-homology pairs (standard
and remote $<$25\% identity) for F1; (ii) the LLaMA-Gene instruction corpus
\citep{wang2024llamagene}, which already follows the candidate-label QA format and
contributes promoter / core-promoter / splice-site / transcription-factor prediction (F3),
protein fold class (7-way, F5), prokaryotic subcellular localization (6-way),
signal-peptide and natural-product-peptide detection (F2), and a Central-Dogma
DNA$\leftrightarrow$protein consistency task (F6); (iii) ProteinGym \citep{notin2023proteingym}
deep-mutational-scanning assays for variant effect (deleterious/benign, F4); (iv)
UniProt-derived protein knowledge QA, Protein2Text sequence-grounded literature QA, and
OPI-Struc sequence$\rightarrow$function generation for F2 function; (v) the bioreason and
protein\_catalogue reasoning corpora for chain-of-thought (F7); (vi) BixBench
\citep{bixbench2024} for biomedical sequence/data-grounded QA (F8); and (vii) the
molecular-image QA from the OmniGene-4-MM multimodal corpus (Vis-CheBI20, F9). Each task
ships train/val/test splits; sequence-level deduplication via identity-aware splitting
guards against train/test leakage. The benchmark is designed for extension: adding a task
requires only a converter that emits schema-valid records and a registry entry, with no
schema change. The full corpus is publicly released as a Hugging Face dataset
(\texttt{dnagpt/biopaws-2}).

\subsection{Dual-mode evaluation protocol}
\label{sec:biopaws2_protocol}

Each model is scored under two modes that, together, give a fair comparison across
paradigms operating each on its home turf.

\textbf{Mode A --- zero-shot QA.} The model answers the test split directly. This measures
instruction-following and innate biological prior. Only generative LLMs can be scored here;
specialized PLMs have no native zero-shot QA interface.

\textbf{Mode B --- fine-tune-then-evaluate.} The model is trained on the task's train split
and evaluated on its test split. A generalist LLM is LoRA-SFT'd on the QA records; a
specialist PLM runs its native protocol (frozen backbone, trained classification/regression
head). Both use the \emph{same} train/test split, each in the training regime it was
designed for.

The leaderboard reports per task the base score (Mode A), the fine-tuned score (Mode B),
and $\Delta = \text{ft} - \text{base}$. The $\Delta$ axis --- the gain a model obtains from
the BioPAWS-2 training data --- is a dimension no prior bio-benchmark reports, and it
directly quantifies \emph{trainability}.

\subsection{The generalist--specialist comparison}
\label{sec:biopaws2_generalist}

The central comparison BioPAWS-2 enables is structural rather than per-cell. A
classification head is a single-task specialist: covering the suite's $N$ tasks requires
$N$ separately trained heads/models. A chat model fine-tuned once on the mixed BioPAWS-2
corpus is a generalist: the same weights answer every task, switching task by instruction
alone. The fair question is therefore not ``does the chat model beat the head on task
$X$?'' but ``is a single generalist model, covering all $N$ tasks with one training run and
one interface, competitive with the ensemble of $N$ task-specific heads?''

Tables~\ref{tab:biopaws2_main} and~\ref{tab:biopaws2_ref} report the result, each under a
single evaluation protocol. A single OmniGene-4-MM, jointly fine-tuned once on the nine-task
mixture and evaluated on the leakage-free splits, attains macro \textbf{0.769} across the
eight discrimination tasks (the free-form F7 chain-of-thought task is kept out of the macro,
as its ROUGE-L is depressed when long-generation competes with short-answer tasks in one
mixture). The headline holds: \emph{one} model, \emph{one} training run, covering the breadth
of the suite. One observation is worth isolating: on the fine-grained biological label sets
(7-way fold class, 6-way localization), \emph{zero-shot} performance is poor for every large
model we tested (Qwen-3.7-max 7.3\% / 23.3\%; OmniGene zero-shot 6.0\% / 50.1\%), because
these specialized label vocabularies are not recoverable from general pretraining --- this is
precisely the gap that the BioPAWS-2 \emph{training corpus} fills, with joint SFT lifting the
generalist to 0.568 / 0.757.

\begin{table}[t]
\centering
\small
\caption{\textbf{BioPAWS-2 joint multi-task SFT (single protocol, leakage-free
entity-disjoint splits).} One OmniGene-4-MM and one raw-Gemma-4-base, each fine-tuned
\emph{once} on the same nine-task mixture and evaluated on every task's clean test split.
Homology/DNA/QA metrics are accuracy; fold/localization are macro-F1; F7 is ROUGE-L. Every
cell is the same protocol (joint SFT) on the same splits --- no per-task or zero-shot
exceptions.}
\label{tab:biopaws2_main}
\begin{tabular}{llcc}
\toprule
Task & Family & OmniGene-4-MM & Gemma-4-base \\
\midrule
Protein homology (standard)      & F1 & \textbf{1.000} & 0.999 \\
Signal peptide                   & F2 & \textbf{0.933} & 0.904 \\
Promoter detection               & F3 & \textbf{0.886} & 0.873 \\
Natural-product peptide          & F2 & \textbf{0.823} & 0.811 \\
Subcellular localization (6-way) & F2 & 0.757 & \textbf{0.775} \\
Central Dogma (DNA$\leftrightarrow$protein) & F6 & 0.604 & 0.604 \\
Variant effect (ProteinGym DMS)  & F4 & 0.578 & \textbf{0.580} \\
Fold class (7-way)               & F5 & \textbf{0.568} & 0.536 \\
\midrule
\textbf{Macro (8 discrimination tasks)} & & \textbf{0.769} & 0.760 \\
\midrule
Bio-CoT reasoning (F7, ROUGE-L)  & F7 & 0.078 & 0.075 \\
\bottomrule
\end{tabular}
\end{table}

\begin{table}[t]
\centering
\small
\caption{\textbf{Reference baselines under their own native protocols.} Specialist ESM-2-3B
head (per-task frozen-backbone + trained head) and Qwen-3.7-max (zero-shot QA), reported
separately because each is a distinct evaluation regime rather than the joint-SFT protocol
of Table~\ref{tab:biopaws2_main}. ESM-2 head values are on the pre-leakage-fix splits and
should be read as optimistic upper bounds; ``---'' = protocol not applicable (ESM-2 is a
protein-only encoder; DNA tasks await a DNA-PLM baseline).}
\label{tab:biopaws2_ref}
\begin{tabular}{llcc}
\toprule
Task & Family & ESM-2-3B head & Qwen-3.7 zero-shot \\
\midrule
Protein homology (standard)      & F1 & 0.998 & 0.813 \\
Protein homology (remote $<$25\%)& F1 & 0.524 & 0.457 \\
Fold class (7-way)               & F5 & \textbf{0.971} & 0.073 \\
Subcellular localization (6-way) & F2 & \textbf{0.954} & 0.233 \\
Signal peptide                   & F2 & \textbf{1.000} & 0.663 \\
Natural-product peptide          & F2 & \textbf{0.929} & --- \\
Variant effect (ProteinGym, acc) & F4 & 0.531 & 0.497 \\
Variant effect (ProteinGym, Spearman) & F4 & 0.36 & --- \\
Promoter detection               & F3 & --- & 0.420 \\
Central Dogma                    & F6 & --- & 0.787 \\
Biomedical QA (BixBench T/F)     & F8 & --- & \textbf{0.937} \\
\bottomrule
\end{tabular}
\end{table}

On protein homology the generalist matches or beats the specialist head; the specialist
remains strongest on structure/localization (fold 0.971, localization 0.954), but each such
number requires a separate trained head, whereas the generalist covers all tasks with one
model. Remote homology is the one task where the generalist's advantage is at
\emph{zero-shot} (0.653, exceeding the specialist head's 0.524 --- see \S\ref{sec:benchmarks})
rather than from SFT; on BioPAWS-2's supervised split, remote-homology SFT collapses to a
majority-class shortcut, consistent with the BioPAWS thesis that the twilight-zone advantage
comes from abductive reasoning over a structural prior, not pattern memorization.

The variant-effect result (F4) is informative about paradigm differences. ProteinGym
deep-mutational-scanning is notoriously hard for mean-pooled PLM embeddings: a single point
mutation barely moves the global mean-pooled ESM-2 representation, so a trained head sits at
chance (0.531). Zero-shot LLMs are also at chance (OmniGene 0.516, Qwen 0.497). Yet a chat
model fine-tuned on the BioPAWS-2 QA records reaches 0.674 ($+16$ points over its zero-shot
baseline): because the mutation is presented as explicit tokens in the instruction rather
than averaged into a pooled vector, the generative paradigm can learn a signal that the
mean-pool-plus-head paradigm structurally discards. This is a concrete case where the
chat-form formulation is not merely a convenience but recovers capability lost by the
classical protocol.

For completeness we also report the field-standard ProteinGym primary metric --- Spearman
correlation between predicted and measured DMS fitness --- rather than only the bucketized
accuracy above. An ESM-2-3B mean-pooled regression head trained on the eight training assays
and evaluated on the held-out assay (AMIE, $n{=}3000$ variants, WT-disjoint from training)
attains Spearman $\rho = 0.36$, a mid-range value typical for single-assay DMS regression at
this model scale. We provide this native-metric number so BioPAWS-2's F4 can be compared
directly against the ProteinGym leaderboard convention; the bucketized-QA form is retained
as the chat-interface-compatible variant, with the raw DMS score preserved per item for
Spearman scoring.

\subsection{Controlling for initialization: is the gain from BioPAWS-2 data or from OmniGene?}
\label{sec:biopaws2_control}

A benchmark that is built and run by the same group that built one of the evaluated models
must guard against a specific confound: when a model performs well after fine-tuning on the
benchmark's training corpus, is the performance attributable to the \emph{model}
(OmniGene's bio-vocabulary and 32.5\,GB biological continued pretraining) or to the
\emph{data} (the BioPAWS-2 SFT corpus itself)? Because the leaderboard reports fine-tuned
scores, this distinction is essential and is exactly the kind of question a fair benchmark
should be able to answer about its own entrants.

We therefore ran a matched control: the \emph{raw} Gemma-4-26B-A4B base model (original
262{,}144 vocabulary, no bio-vocabulary injection, no biological continued pretraining) was
fine-tuned with the \emph{identical} joint multi-task LoRA recipe (same nine-task mixture,
same per-task cap, same $r$=64/$\alpha$=128 adapter, same learning rate, epochs, and
batch size) as OmniGene-4-MM. The only differences are the initialization and the
tokenizer. Table~\ref{tab:biopaws2_control} reports the result.

\begin{table}[t]
\centering
\small
\caption{\textbf{Initialization control under identical joint SFT.} OmniGene-4-MM
(bio-vocab + 32.5\,GB bio-CPT) vs.\ raw Gemma-4 base (no bio-CPT), both LoRA-fine-tuned on
the same nine-task BioPAWS-2 mixture with the same recipe, on the leakage-free
entity-disjoint splits (\S\ref{sec:biopaws2_leakage}). Macro is over the eight
discrimination tasks (the free-form F7 chain-of-thought task is excluded from the macro;
see text). The same $\approx$-zero gap held on the pre-leakage-fix splits (macro
0.783 vs.\ 0.772, $\Delta{+}0.011$), so the conclusion is robust to the split.}
\label{tab:biopaws2_control}
\begin{tabular}{lccc}
\toprule
Task & OmniGene-4-MM + SFT & Gemma-4-base + SFT & $\Delta$ \\
\midrule
Protein homology (standard)  & 1.000 & 0.999 & $+0.001$ \\
Subcellular localization     & 0.757 & 0.775 & $-0.018$ \\
Promoter detection           & 0.886 & 0.873 & $+0.013$ \\
Fold class (7-way)           & 0.568 & 0.536 & $+0.032$ \\
Central Dogma                & 0.604 & 0.604 & $+0.000$ \\
Signal peptide               & 0.933 & 0.904 & $+0.030$ \\
Natural-product peptide      & 0.823 & 0.811 & $+0.012$ \\
Variant effect (ProteinGym)  & 0.578 & 0.580 & $-0.002$ \\
\midrule
\textbf{Macro (8 tasks)}     & \textbf{0.769} & \textbf{0.760} & $+0.008$ \\
\bottomrule
\end{tabular}
\end{table}

The two models are \emph{statistically indistinguishable} under joint SFT: the macro gap is
$+0.008$, the mean per-task difference is $+0.008$, the largest single-task difference is
$0.032$ (fold class), and the sign of the difference is inconsistent across tasks (five
favor OmniGene, two are ties, two favor the base model). We read this honestly: \textbf{on
BioPAWS-2's supervised tasks, almost all of the achievable performance comes from the
BioPAWS-2 training data, not from OmniGene's biological pretraining}. Biological CPT does
not provide a material advantage \emph{once a model is fine-tuned on the instruction
corpus}.

This is a useful negative result that the benchmark itself surfaces, and it sharpens --- rather
than weakens --- the contribution of this paper. First, it cleanly separates the two claims:
BioPAWS-2 is established as an effective, model-agnostic training resource (a strong base
model with no domain pretraining reaches macro 0.760), while OmniGene-4-MM's distinctive
value is shown to lie \emph{elsewhere} than in supervised-fine-tuning headroom --- namely in
its \emph{zero-shot} behaviour (e.g.\ remote-homology detection, \S\ref{sec:benchmarks},
where it exceeds specialist PLMs without any task training) and in its router-level
interpretability (\S\ref{sec:moe}), neither of which the base model possesses. Second, it
directly answers the self-benchmarking concern with data rather than rhetoric: the group
that built the benchmark used it to show that its \emph{own} model's pretraining confers no
SFT advantage on these tasks --- the opposite of a benchmark tuned to flatter its authors.

\subsection{Positioning}

BioPAWS-2 unifies the task content of TAPE, PEER, FLIP, ProteinGym, GUE, DeepLoc, SignalP,
and the LLaMA-Gene corpus into a single chat-formatted axis on which a specialist PLM (in
its native head-training regime) and a generalist LLM (zero-shot or fine-tuned) can be
compared directly, with a dual zero-shot/fine-tune protocol and an explicit
initialization control (\S\ref{sec:biopaws2_control}). To our knowledge it is the first
biological benchmark to combine instruction-tuning-native format, SFT-trainability,
multimodality, and cross-modality in one suite, though we make this breadth claim
descriptively rather than as a priority claim. By releasing it as a public training
resource on which any model can be fine-tuned and scored --- and by using it to report a
negative control on our own model's pretraining --- BioPAWS-2 makes OmniGene-4-MM one
entrant among many on an open leaderboard rather than the author of its own test, which is
how we address the self-benchmarking concern of the previous submission.
